
% ============================================================================
\clearpage
\section*{v4.8.0-r08 Source-Only: Test of Prediction P1 (Photon-Frame Unfolding)}
\addcontentsline{toc}{section}{v4.8.0-r08 Source-Only: Test of Prediction P1 (Photon-Frame Unfolding)}
\label{sec:v480-r08-p1-test}
% ============================================================================

\paragraph{Purpose.}
Corollary~\ref{cor:falsifiable} lists P1 as testable now: unfolding the zeta zeros with the
photon-frame density should improve the GUE fit by about \(21\%\) (Corollary~\ref{cor:odlyzko}),
and Theorem~\ref{thm:gue-selfdet} states that the Kolmogorov--Smirnov distance has a unique global
minimum at \(\mu=\gL\).  This layer performs the test on the first 2000 zeros.  Earlier layers are
not rewritten.

\noindent\textbf{Source-only status.}
This v4.8.0-r08 layer is source-only.  It does not build a PDF, does not assign a new DOI,
and does not authorize a public release.  The v4.8.0 Version DOI
\href{https://doi.org/10.5281/zenodo.23365244}{10.5281/zenodo.23365244} and the Concept DOI
\href{https://doi.org/10.5281/zenodo.19210728}{10.5281/zenodo.19210728} are unchanged by it.

\subsection*{v4.8.0-r08 setup}
\label{subsec:v480-r08-setup}

As in Theorem~\ref{thm:gue-selfdet},
\(\tilde\gamma_n(\mu)=\int_0^{\gamma_n}\frac{\log(t/2\pi\mu)}{2\pi\mu}\,dt
=\frac{\gamma_n}{2\pi\mu}\bigl(\log\frac{\gamma_n}{2\pi\mu}-1\bigr)\) and
\(s_n(\mu)=\tilde\gamma_{n+1}(\mu)-\tilde\gamma_n(\mu)\).  The comparison law is the Wigner surmise
\(p(s)=\frac{32}{\pi^2}s^2e^{-4s^2/\pi}\); \(\chi^2\)/dof uses 15 bins on \([0,3]\).  The scale runs over
\(\mu\in[0.5,3]\) in steps of \(0.05\).  Two readings are tested: (a) raw, exactly as defined, and
(b) normalized to mean spacing \(1\), since a pure rescaling would otherwise move the mean.

\subsection*{v4.8.0-r08 result}
\label{subsec:v480-r08-result}

\begin{center}
\small
\begin{tabular}{@{}>{\raggedright\arraybackslash}p{0.26\textwidth}>{\raggedright\arraybackslash}p{0.33\textwidth}>{\raggedright\arraybackslash}p{0.33\textwidth}@{}}
\toprule
 & zeros 1--100 & zeros 101--2000\\
\midrule
(a) raw: KS minimum & \(\mu=1.00\) & \(\mu=1.00\)\\
(a) KS at \(\mu=1\) / \(\mu=\gL\) & \(0.068\) / \(0.699\) & \(0.041\) / \(0.628\)\\
(a) mean spacing at \(\mu=\gL\) & \(0.419\) & \(0.473\)\\
(b) normalized: KS minimum & \(\mu=2.85\) & \(\mu=2.25\)\\
(b) KS range over \(\mu\) & \(0.037\)--\(0.070\) & \(0.039\)--\(0.042\)\\
(b) \(\chi^2\)/dof change \(\mu=1\to\gL\) & \(+5.4\%\) improvement & \(-9.9\%\) (worse)\\
\bottomrule
\end{tabular}
\end{center}

\paragraph{Theoremlet v4.8.0-r08.1.}
On the first 2000 zeros, neither reading of the photon-frame unfolding singles out \(\mu=\gL\).
In reading (a) the unique minimum is the standard unfolding \(\mu=1\); in reading (b) the
dependence on \(\mu\) is flat within the statistical scatter and its minimum is not at \(\gL\).
Moreover, in reading (b) the effect of \(\mu\) vanishes with height: the local stretch it introduces
is of relative size \(\log\mu/\log(\gamma/2\pi)\), about \(1\%\) at \(\gamma\sim10^{22}\), so an
asymptotic improvement of \(21\%\) cannot occur.

\paragraph{Proof sketch.}
Direct computation, script \texttt{k11\_p1\_entfaltung.py} in the working folder.  For reading (b),
\(\rho_\mu(t)/\rho_1(t)=\mu^{-1}\bigl(1-\log\mu/\log(t/2\pi)\bigr)\); after normalization the constant
factor \(\mu^{-1}\) drops out and only the second term, which decays like \(1/\log t\), remains.
\(\square\)

\paragraph{Reproducibility note.}
The value of \(75.8\%\) for 100 zeros stated in Corollary~\ref{cor:odlyzko} is not reproduced in
either reading.  The original script \texttt{gue\_frame\_correction.py} is not in the source folders,
so its exact procedure could not be compared.

\subsection*{v4.8.0-r08 consequences}
\label{subsec:v480-r08-consequences}

\begin{center}
\small
\begin{tabular}{@{}>{\raggedright\arraybackslash}p{0.34\textwidth}>{\raggedright\arraybackslash}p{0.58\textwidth}@{}}
\toprule
statement & status after r08\\
\midrule
Thm.~\ref{thm:gue-selfdet} (unique minimum at \(\gL\)) & not confirmed; contradicted on zeros 1--2000\\
Cor.~\ref{cor:odlyzko}, P1 in Cor.~\ref{cor:falsifiable} & by its own criterion falsified at low height; the asymptotic form excluded by the \(1/\log\gamma\) decay\\
\(\gL\) as ``GUE fixpoint of the frame'' & withdrawn; \(\gL\) remains an input (cf.\ r03 F11, OP1)\\
\bottomrule
\end{tabular}
\end{center}

Corollary~\ref{cor:odlyzko} states that non-observation falsifies \(\gL=\gBI/\gstring\) as the
physical frame factor.  This layer records the non-observation on the accessible zeros and the
reason why the stated asymptotic effect cannot appear.  It does not by itself decide the physical
status of \(\gL\); that remains with OP1.

\subsection*{v4.8.0-r08 ledger and guard update}
\label{subsec:v480-r08-ledger-guard}

\[
R_{\rm P1}^{v480\text{-}r08}=(0,0,1)
\]
for (test performed; prediction not confirmed; physical status of \(\gL\) open).

The layer spends no guard:
\[
R_{\rm glob}^{(v480\text{-}r08,+)}=(0,1,1,1,1),
\qquad
\Gamma_{\rm guard}^{v480\text{-}r08}=(1,1,1,1,1),
\qquad
\beta_4=1/256\ \text{unspent}.
\]
It does not imply exact finite \(C5\) zero closure, exact all-readout terminality,
\(C6/T2\), a global photon-clock/RH theorem, or global all-readout closure.

